\section{Introduction}

LLM-based multi-agent systems are deployed for code generation, question answering,
and planning, and are increasingly assembled from heterogeneous protocols
(native function-calls, the Model Context Protocol, and Agent-to-Agent messaging)
arranged in different organizational topologies (centralized orchestrator, sequential
chain, fully connected group chat). Recent work\,\cite{scaling_agents_2025} reports
aggregate failure rates of 41--86\% across multi-agent benchmarks, but the
contribution of \emph{communication protocol} versus \emph{organizational topology}
to these failures is not separately quantified.

We disentangle these two design axes. Holding the underlying LLM, agent count, and
task suite fixed, we run a $3 \times 3$ factorial sweep over protocol
$\in \{\textsc{native}, \textsc{mcp}, \textsc{a2a}\}$ and topology
$\in \{\textsc{centralized}, \textsc{chain}, \textsc{fully\_connected}\}$, and
classify every trial trace with the MAST taxonomy
(FC1: specification issues, FC2: inter-agent misalignment, FC3: task verification
failures). Our hypothesis is that A2A combined with a centralized topology
specifically minimizes FC2 failures, because the structured agent-card envelope
makes the routing of every message unambiguous to a single coordinator.

Our contributions are: (1) a self-contained, no-autogen replication harness for
protocol $\times$ topology MAST experiments; (2) a $3 \times 3$ FC2 heatmap with
per-cell FC1/FC3 breakdowns; (3) two follow-up sweeps over agent count
($N \in \{2,3,4,5\}$) and task type ($\{\textsc{code}, \textsc{qa},
\textsc{planning}\}$) that test the generality of the best cell.
